\begin{center}
{\small 公司内部发表前审批用}\par\vspace{7pt}
{\LARGE\bfseries 空中-水面-水下无人系统中的区块链辅助任务信任：\\[3pt]综述与概念框架}\par\vspace{8pt}
{\large 中文全文翻译}\par\vspace{9pt}
Shaowei Liu\textsuperscript{1,2}，Jing Wu\textsuperscript{1}，Tanren Liu\textsuperscript{3}，Shengqiang He\textsuperscript{3}，Yakun Ren\textsuperscript{3}，Chengnian Long\textsuperscript{1,*}
\end{center}
\begingroup\small\noindent
\textsuperscript{1} 上海交通大学，School of Automation and Intelligent Sensing，中国上海市东川路 800 号，200240；liusw1994@sjtu.edu.cn（S.L.）；jingwu@sjtu.edu.cn（J.W.）\\
\textsuperscript{2} 江苏自动化研究所，中国江苏省连云港市圣湖路 18 号，222061。\\
\textsuperscript{3} 顺丰科技有限公司（SF Technology Co., Ltd.），中国广东省深圳市南山区科技南一路 TK Chuangzhi Tiandi 大厦 B 座，518057；tanrenliu@sf-express.com（T.L.）；heshengqiang@sf-express.com（S.H.）；yakunren@sf-express.com（Y.R.）。\\
* 通讯作者：longcn@sjtu.edu.cn。\par
\medskip\noindent 学术编辑：Maurizio Talamo。收稿：2026 年 9 月 30 日；修回：2026 年 10 月 28 日；接受：2026 年 10 月 31 日。\par
\medskip\noindent 译稿说明：本译稿对应前附的 2026 年 11 月 2 日作者校正英文稿，供公司内部发表前审批。章节、公式、图表及引文编号与英文稿对应；图片保留英文原图，图题译为中文；人名、标准编号、软件名称及参考文献原始书目信息保留，参考文献另附中文题名。\par\endgroup
\vspace{8pt}
\noindent\textbf{摘要：}空中-水面-水下任务将异构无人平台、间歇性链路和多个管理域连接起来。可靠的决策必须将物理观测、受保护的记录、共享权限以及用途适合性联系起来。本结构化叙述性综述通过 89 条限定范围的主张记录，将对象-保障-生命周期（Object--Assurance--Lifecycle，OAL）框架应用于 39 项经过全文检查的机制研究。比较性综合区分了发送者认证与观测有效性、平台完整性与任务适用性，以及复制状态恢复与任务数据失效处理。按领域分层的分析和敏感性分析揭示了纳入选择与编码选择对覆盖图谱的影响。依赖图识别受输入或权限变化影响的数据产品，情境-时间-图（CTG）循环则组织对这些产品的重新评估。分区-撤销-协调案例明确了联盟角色、面向 Fabric 的应用模式、资源受限条件下的验证，以及声学链路的载荷预算。合成事件回放验证了签名、包含证明、时间规则和授权规则。在规则一致的条件下，具备策略验证能力的见证日志产生了相同的决策。因此，账本的价值取决于共同治理需求及实现成本。任务的故障与通信要求指导共识选择；提出的测试在匹配条件下比较决策质量与恢复能力。本综述将机制证据与跨域任务决策的明确组合要求联系起来。

\noindent\textbf{关键词：}任务信任；无人系统；UAV-USV-UUV 协作；区块链；数据溯源；远程证明；零信任；海事自主系统
